% Wersja źródłowa: audyt II, jednostka 33
\begingroup
\setlength{\parskip}{3pt}
\setcounter{section}{32}
\noindent\textbf{Wyniki wykorzystywane.} Z poprzednich rozdziałów korzystamy z twierdzenia o wykresie domkniętym, teorii operatorów sprzężonych na przestrzeniach Hilberta, twierdzenia Plancherela i twierdzenia spektralnego dla operatorów ograniczonych normalnych wraz z rachunkiem borelowskim. Dowody twierdzenia spektralnego dla operatorów nieograniczonych, twierdzenia von Neumanna o rozszerzeniach i twierdzenia Stone'a podajemy w tym rozdziale.\par\medskip
\section{Operatory nieograniczone na przestrzeniach Hilberta}
\label{sec:unbounded33}

W \hyperref[sec:spectral-hilbert32]{rozdziale~32} skonstruowaliśmy miarę spektralną
operatora \emph{ograniczonego} i normalnego. Operator różniczkowania na
$L^2$, operator mnożenia przez zmienną oraz generatory jednoparametrowych
grup unitarnych są na ogół nieograniczone. Zasadnicza różnica polega na tym,
że \emph{dziedzina operatora jest częścią jego definicji}. W całym rozdziale
$H$ oznacza niezerową zespoloną przestrzeń Hilberta, a iloczyn skalarny jest liniowy
w pierwszym argumencie. Zapis $\mathcal L(H)$ oznacza operatory ograniczone
na całym $H$.

\subsection{Dziedziny, wykresy, operatory sprzężone}

\begin{definicja}[Operator liniowy i jego wykres]
\label{def:unbounded-operator33}
Operatorem liniowym w $H$ nazywamy liniowe odwzorowanie
$A:D(A)\to H$, gdzie $D(A)\subseteq H$ jest podprzestrzenią liniową.
Jego \emph{wykresem} jest $G(A)=\{(x,Ax):x\in D(A)\}\subseteq H\times H$.
Operator jest \emph{domknięty}, gdy $G(A)$ jest domknięty, oraz
\emph{domykalny}, gdy domknięcie $G(A)$ jest wykresem pewnego operatora;
operator ten nazywamy \emph{domknięciem} $\overline A$.
Dla operatorów $A,B$ piszemy $A\subseteq B$, jeżeli
$D(A)\subseteq D(B)$ i $Bx=Ax$ dla $x\in D(A)$.
\end{definicja}

\begin{stwierdzenie}[Kryteria domkniętości i norma wykresowa]
\label{prop:graph33}
Niech $A:D(A)\subseteq H\to H$ będzie liniowy.
Wówczas $A$ jest domknięty wtedy i tylko wtedy, gdy dla każdego
ciągu $x_n\in D(A)$ z $x_n\to x$ i $Ax_n\to y$ w $H$ zachodzi
$x\in D(A)$ i $Ax=y$. Ponadto $A$ jest domknięty wtedy i tylko wtedy,
gdy $D(A)$ z normą $\|x\|_A=(\|x\|^2+\|Ax\|^2)^{1/2}$ jest
przestrzenią Hilberta z iloczynem skalarnym
$\langle x,y\rangle_A=\langle x,y\rangle+\langle Ax,Ay\rangle$.
Jeżeli $D(A)=H$, to operator domknięty jest ograniczony.
\end{stwierdzenie}
\begin{proof}
Pierwsza równoważność jest ciągowym kryterium domkniętości
podprzestrzeni przestrzeni metrycznej $H\times H$.
Odwzorowanie $D(A)\ni x\mapsto(x,Ax)\in G(A)$ jest izometrią
przy wskazanych iloczynach skalarnych, a $G(A)$ jest zupełny dokładnie
wtedy, gdy jest domknięty. Ostatnia teza wynika z twierdzenia
o wykresie domkniętym dla przestrzeni Banacha
(\hyperref[sec:frechet-open-graph]{rozdział~11}).
\end{proof}

\begin{definicja}[Zbiór rezolwenty operatora domkniętego]
Niech $A:D(A)\subseteq H\to H$ będzie operatorem domkniętym.
Do jego \emph{zbioru rezolwenty} $\rho(A)$ należą te $z\in\C$,
dla których $A-zI:D(A)\to H$ jest bijekcją, a jej
odwrotność $(A-zI)^{-1}:H\to H$ jest ograniczona.
\emph{Widmo} to $\sigma(A)=\C\setminus\rho(A)$.
Dla operatora domkniętego ograniczoność odwrotności wynika już
z jej istnienia i twierdzenia o wykresie domkniętym,
ale podajemy ją jawnie dla odróżnienia od operatorów niedomkniętych.
\end{definicja}

\begin{definicja}[Operator sprzężony]
\label{def:adjoint33}
Niech $A:D(A)\subseteq H\to H$ będzie liniowy i gęsto określony,
czyli $\overline{D(A)}=H$. Definiujemy
\[
 D(A^*)=\{y\in H:\text{istnieje }z\in H\text{ takie, że }
 \langle Ax,y\rangle=\langle x,z\rangle\ (x\in D(A))\},
 \qquad A^*y=z.
\]
Wektor $z$ jest jednoznaczny na mocy gęstości $D(A)$.
Operator $A$ jest \emph{symetryczny}, gdy $A\subseteq A^*$,
a \emph{samosprzężony}, gdy $A=A^*$, włącznie z równością dziedzin.
\end{definicja}

\begin{twierdzenie}[Sprzężenie i domykalność]
\label{thm:closable-adjoint33}
Niech $A:D(A)\subseteq H\to H$ będzie gęsto określony.
Wówczas $A^*$ jest operatorem domkniętym (jego dziedzina
nie musi być gęsta). Operator $A$ jest domykalny wtedy i tylko wtedy,
gdy $D(A^*)$ jest gęsta w $H$. W tym przypadku
$\overline A=A^{**}$. Każdy gęsto określony operator symetryczny
jest domykalny, a jego domknięcie również jest symetryczne.
\end{twierdzenie}
\begin{proof}
Jeżeli $y_n\to y$ i $A^*y_n\to z$, to dla każdego $x\in D(A)$
$\langle Ax,y\rangle=\lim_n\langle Ax,y_n\rangle
=\lim_n\langle x,A^*y_n\rangle=\langle x,z\rangle$;
stąd $y\in D(A^*)$ i $A^*y=z$.

Sprawdźmy kryterium domykalności. Dla wektora $(u,v)\in H\times H$
warunek $(u,v)\perp G(A)$ jest równoważny
$\langle u,x\rangle+\langle v,Ax\rangle=0$ dla wszystkich
$x\in D(A)$. Po sprzężeniu tej równości otrzymujemy
$v\in D(A^*)$ oraz $u=-A^*v$. Zatem
$G(A)^\perp=\{(-A^*v,v):v\in D(A^*)\}$.
W szczególności $(0,w)\in\overline{G(A)}$ wtedy i tylko wtedy,
gdy $\langle w,v\rangle=0$ dla wszystkich $v\in D(A^*)$.
Liniowa podprzestrzeń $\overline{G(A)}$ jest wykresem operatora
dokładnie wtedy, gdy nie zawiera niezerowych wektorów $(0,w)$.
Otrzymujemy równoważność z gęstością $D(A^*)$.
Gdy ta gęstość zachodzi, dla $(x,y)\in H\times H$ mamy
\[
 (x,y)\in G(A)^{\perp\perp}
 \ \Longleftrightarrow\
 \langle y,v\rangle=\langle x,A^*v\rangle
 \quad(v\in D(A^*)).
\]
Po sprzężeniu jest to dokładnie warunek $x\in D(A^{**})$
i $A^{**}x=y$. Zatem $G(A^{**})=\overline{G(A)}$.
Jeżeli $A$ jest symetryczny, to $D(A)\subseteq D(A^*)$, a więc
$D(A^*)$ jest gęsta; przejście do granic w tożsamości
$\langle Ax,y\rangle=\langle x,Ay\rangle$ zachowuje symetrię.
\end{proof}

\begin{uwaga}[Operator ograniczony na niewłaściwej dziedzinie]
Nawet operator zerowy na gęstej właściwej podprzestrzeni $D\subsetneq H$
nie jest samosprzężony: jego operator sprzężony jest zerowy na całym
$H$, natomiast jego domknięcie jest operatorem zerowym na $H$.
Nie wystarczy zatem sprawdzać samego wzoru działania operatora.
\end{uwaga}

\subsection{Symetria, samosprzężoność i rezolwenta}

\begin{lemat}[Oszacowanie poza osią rzeczywistą]
\label{lem:symmetric-resolvent33}
Niech $A$ będzie gęsto określonym operatorem symetrycznym w $H$.
Dla $z=a+ib\in\C$ z $b\ne0$ i $x\in D(A)$ zachodzi
$\|(A-zI)x\|\ge |b|\|x\|$. Jeśli $A$ jest ponadto domknięty,
$\operatorname{Ran}(A-zI)$ jest podprzestrzenią domkniętą.
\end{lemat}
\begin{proof}
Ponieważ $\langle Ax,x\rangle$ jest rzeczywiste,
$|\operatorname{Im}\langle(A-zI)x,x\rangle|=|b|\|x\|^2$.
Nierówność Cauchy'ego--Schwarza daje oszacowanie.
Gdy $A$ jest domknięty, oszacowanie implikuje, że
$(A-zI)x_n\to y$ pociąga zbieżność $x_n\to x$;
ponadto $Ax_n=(A-zI)x_n+zx_n\to y+zx$, więc
domkniętość $A$ daje $x\in D(A)$ i $(A-zI)x=y$.
\end{proof}

\begin{twierdzenie}[Kryterium samosprzężoności]
\label{thm:selfadjoint-range33}
Niech $A$ będzie gęsto określonym, symetrycznym operatorem w $H$.
Wówczas następujące warunki są równoważne:
\begin{enumerate}[label=\textup{(\roman*)}]
\item $A$ jest samosprzężony;
\item $\operatorname{Ran}(A+iI)=H$ i $\operatorname{Ran}(A-iI)=H$.
\end{enumerate}
Jeśli warunki te zachodzą, dla każdego $z\in\C\setminus\R$
operator $A-zI:D(A)\to H$ jest bijekcją, a jego odwrotność
$R(z,A)=(A-zI)^{-1}\in\mathcal L(H)$ spełnia
$\|R(z,A)\|\le1/|\operatorname{Im}z|$.
\end{twierdzenie}
\begin{proof}
Jeżeli $A=A^*$, to $A$ jest domknięty. Z lematu
\ref{lem:symmetric-resolvent33} oba obrazy $A\pm iI$ są domknięte.
Ogólnie, dla $z\in\C$, z definicji sprzężenia wynika
\[
 \operatorname{Ran}(A-zI)^\perp=\ker(A^*-\overline z I).
\]
Istotnie, $\langle (A-zI)x,y\rangle=0$ dla każdego $x\in D(A)$
jest równoważne $\langle Ax,y\rangle=\langle x,\overline z y\rangle$.
Ortogonalne dopełnienia obrazów $A\pm iI$ wynoszą więc odpowiednio
$\ker(A^*\mp iI)$, a te są zerowe: dla $y\in D(A)$
równość $(A\mp iI)y=0$ i oszacowanie z lematu wymuszają $y=0$.
Obrazy są więc również gęste, a zatem równe $H$.

Odwrotnie, załóżmy surjektywność obu operatorów.
Wówczas $A$ jest domknięty: jeśli $x_n\to x$ i $Ax_n\to y$,
to $(A+iI)x_n\to y+ix$; wybierając $u\in D(A)$
z $(A+iI)u=y+ix$ i stosując oszacowanie z lematu,
otrzymujemy $x_n\to u$, więc $x=u\in D(A)$ i $Ax=y$.
Dla $y\in D(A^*)$ wybieramy $u\in D(A)$ z
$(A+iI)u=(A^*+iI)y$. Wówczas
$(A^*+iI)(y-u)=0$, a
$\ker(A^*+iI)=\operatorname{Ran}(A-iI)^\perp=\{0\}$.
Stąd $y=u\in D(A)$ i $A=A^*$.

Dla dowolnego $z\notin\R$ ponawiamy argument z domkniętością
obrazu i jego gęstością, używając $A=A^*$ oraz oszacowania
z lematu. Oszacowanie normy odwrotności jest natychmiastowe.
\end{proof}

\begin{przyklad}[Mnożenie przez funkcję rzeczywistą]
\label{ex:mult-unbounded33}
Niech $(X,\Sigma,\mu)$ będzie przestrzenią z miarą, a
$m:X\to\R$ funkcją mierzalną, skończoną prawie wszędzie.
W przestrzeni Hilberta $H=L^2(X,\mu)$ operator
$(M_mf)(x)=m(x)f(x)$ z dziedziną
$D(M_m)=\{f\in L^2(X,\mu):mf\in L^2(X,\mu)\}$
jest gęsto określony i samosprzężony. Istotnie, funkcje
$f\chi_{\{|m|\le n\}}$ należą do $D(M_m)$ i zbiegają do $f$
w $L^2$. Dla każdego $g\in H$ funkcje
$f_\pm(x)=g(x)/(m(x)\pm i)$ należą do $D(M_m)$,
ponieważ $|m/(m\pm i)|\le1$, i spełniają
$(M_m\pm iI)f_\pm=g$. Operator jest symetryczny, więc
poprzednie twierdzenie daje samosprzężoność.
Dla $X=\R$ z miarą Lebesgue'a $\lambda_1$ i $m(x)=x$
operator jest nieograniczony: wektory $f_n=\chi_{[n,n+1]}$
należą do dziedziny, mają normę $1$, a $\|M_m f_n\|\ge n$.
\end{przyklad}

\subsection{Miary spektralne i rachunek funkcji nieograniczonych}

\begin{lemat}[Całkowanie nieograniczonych funkcji względem miary spektralnej]
\label{lem:unbounded-integral33}
Niech $K\subseteq\C$ będzie niepustym zbiorem borelowskim
(niekoniecznie zwartym), $H$ zespoloną przestrzenią Hilberta,
a $E:\mathcal B(K)\to\mathcal L(H)$ miarą spektralną:
$E(\varnothing)=0$, $E(K)=I_H$, każde $E(B)$ jest rzutem ortogonalnym,
$E(B\cap C)=E(B)E(C)$ i $E$ jest przeliczalnie addytywna
w silnej topologii operatorowej. Niech $f:K\to\C$ będzie
funkcją borelowską, niekoniecznie ograniczoną.
Wówczas istnieje gęsto określony operator domknięty $f(E)$
o dziedzinie
\[
 D(f(E))=\left\{x\in H:\int_K|f(z)|^2\,\mathrm d\mu_x(z)<\infty\right\}.
\]
Dla $x\in D(f(E))$ zachodzi
\[
 \|f(E)x\|^2=\int_K|f(z)|^2\,\mathrm d\mu_x(z),
\]
gdzie $\mu_x(B)=\langle E(B)x,x\rangle$.
Zachodzi $f(E)^*=\overline f(E)$; w szczególności dla
$f$ rzeczywistej $E$-prawie wszędzie operator $f(E)$ jest
samosprzężony. Dla ograniczonej borelowskiej funkcji $g$
operator $g(E)$ zachowuje $D(f(E))$ oraz
$f(E)g(E)=(fg)(E)$ z równością dziedzin. Ponadto
$g(E)f(E)x=(fg)(E)x$ dla $x\in D(f(E))$.
Dla ograniczonej $f$ konstrukcja jest zwykłym borelowskim
rachunkiem funkcyjnym.
\end{lemat}
\begin{proof}
\textbf{Krok 1: rachunek funkcji ograniczonych.}
Dla funkcji prostej $s=\sum_{j=1}^m c_j\chi_{B_j}$, gdzie
$B_j$ są parami rozłączne, definiujemy
$s(E)=\sum_{j=1}^m c_jE(B_j)$.
Dopełniając rozłączny układ do podziału $K$ (ze współczynnikiem
zero na pozostałej części) i stosując wspólne rozdrobnienie
dwóch podziałów, otrzymujemy niezależność od reprezentacji
funkcji prostej oraz liniowość. Argument jest taki sam jak
w Lemacie~\ref{lem:pvm-integral32}; nie wymaga zwartości $K$.
Ponadto $1(E)=I_H$. Własności rzutów dają
$\|s(E)x\|^2=\int_K|s|^2\,\mathrm d\mu_x$,
$(st)(E)=s(E)t(E)$ i $s(E)^*=\overline s(E)$
dla funkcji prostych. Każdą ograniczoną funkcję borelowską
można jednostajnie aproksymować funkcjami prostymi.
Z nierówności $\|s(E)\|\le\|s\|_\infty$ otrzymujemy
przez granicę w normie operatorowej rachunek wszystkich
ograniczonych funkcji borelowskich, z tymi samymi
własnościami algebraicznymi i normowymi.

\textbf{Krok 2: dziedzina i domkniętość.}
Funkcja $f$ jest skończona punktowo. Dla ograniczonych
funkcji równych poza zbiorem $N$ z $E(N)=0$ równość operatorów
wynika ze wzoru na normę ich różnicy, ponieważ $\mu_x(N)=0$
dla każdego $x\in H$. Po modyfikacji dowolnej funkcji $f$ na
takim zbiorze nie zmieniają się całki $|f|^2$ ani operatory
jej ograniczonych obcięć, a więc również poniższa konstrukcja
dziedziny i granicy.
Kładziemy $B_n=\{z\in K:|f(z)|\le n\}$,
$P_n=E(B_n)$ i $f_n=f\chi_{B_n}$.
Wtedy $P_n\to I_H$ silnie, a $f_n$ są ograniczone.
Jeżeli $x$ należy do przestrzeni $D$ określonej całką w tezie,
to
\[
 \|(f_n(E)-f_m(E))x\|^2
 =\int_K|f_n-f_m|^2\,\mathrm d\mu_x\longrightarrow0.
\]
Definiujemy $f(E)x=\lim_n f_n(E)x$.
Odwrotnie, gdy $f_n(E)x$ ma granicę, to jego normy są
ograniczone, a monotoniczna zbieżność całek
$\int|f_n|^2\,\mathrm d\mu_x$ dowodzi $x\in D$.
Ponieważ $P_nH\subseteq D$ (dla $x\in H$ miara $\mu_{P_nx}$
jest skupiona na $B_n$ i $\int|f|^2\,\mathrm d\mu_{P_nx}\le n^2\|x\|^2$)
i $P_n\to I_H$ silnie, dziedzina $D$ jest gęsta.
Jest ona podprzestrzenią liniową, gdyż równoważnie składa się
z tych $x$, dla których ciąg liniowych odwzorowań $f_n(E)x$
jest zbieżny. Operator graniczny jest więc liniowy.
Tożsamość na normę wynika
z monotonicznej zbieżności całek.

Aby wykazać domkniętość, przypuśćmy, że $x_j\in D$,
$x_j\to x$ oraz $f(E)x_j\to y$. Dla każdego $n$
z rachunku ograniczonego wynika
$P_nf(E)x_j=f_n(E)x_j$.
Przechodząc do granicy przy ustalonym $n$ otrzymujemy
$P_ny=f_n(E)x$, a więc $\|f_n(E)x\|\le\|y\|$.
Monotoniczna zbieżność całek pokazuje, że $x\in D$.
Następnie $f_n(E)x=P_ny\to y$, czyli $f(E)x=y$.

\textbf{Krok 3: operator sprzężony i mnożenie.}
Dla $x\in D(f(E))$ i $y\in D(\overline f(E))$,
przechodząc do granicy w tożsamości dla $f_n(E)$,
dostajemy $\langle f(E)x,y\rangle
=\langle x,\overline f(E)y\rangle$.
Zatem $\overline f(E)\subseteq f(E)^*$.
Jeśli zaś $y\in D(f(E)^*)$ i $z=f(E)^*y$, to
testując tożsamość sprzężenia na $P_nx$ dla dowolnego
$x\in H$ otrzymujemy
$\overline{f_n}(E)y=P_nz$.
Stąd $\|\overline{f_n}(E)y\|\le\|z\|$.
Monotoniczna zbieżność całek implikuje
$y\in D(\overline f(E))$, a granica daje
$\overline f(E)y=z$. Zatem $f(E)^*=\overline f(E)$.
Jeżeli $g$ jest ograniczona, to z rachunku ograniczonego
\[
 \mu_{g(E)x}(C)=\|(\chi_Cg)(E)x\|^2
 =\int_C |g|^2\,\mathrm d\mu_x
 \quad(C\in\mathcal B(K)).
\]
Stąd $\int|f|^2\,\mathrm d\mu_{g(E)x}=\int|fg|^2\,\mathrm d\mu_x$
(przez aproksymację funkcjami prostymi i zbieżność monotoniczną).
W szczególności $D(f(E)g(E))=D((fg)(E))$, a dziedzina ta zawiera
$D(f(E))$, ponieważ prawa całka jest wtedy nie większa niż
$\|g\|_\infty^2\int|f|^2\,\mathrm d\mu_x$.
Uzasadnijmy przy tym ogólną zasadę obcięć: jeżeli $h$ jest
borelowska, $x\in D(h(E))$, a ograniczone funkcje $h_n$
zbiegają do $h$ punktowo i spełniają $|h_n|\le |h|$, to
\[
 \|h_n(E)x-h(E)x\|^2=\int_K|h_n-h|^2\,\mathrm d\mu_x.
\]
Istotnie, dla ustalonego $n$ stosujemy wzór dla funkcji
ograniczonych $h_n$ oraz $h\chi_{\{|h|\le m\}}$ i przechodzimy
z $m$ do nieskończoności. Po stronie operatorowej korzystamy
z definicji $h(E)$, a po stronie całkowej ze zbieżności
zmajoryzowanej przez $4|h|^2$. Kolejne przejście z $n$ do
nieskończoności daje zbieżność w normie $H$.
Stosując tę zasadę do $h=fg$ i $h_n=f_ng$, otrzymujemy
$(f_ng)(E)x\to(fg)(E)x$ dla $x\in D((fg)(E))$. Z drugiej strony $(f_ng)(E)x=f_n(E)g(E)x$,
co dowodzi równości $f(E)g(E)=(fg)(E)$ z równością dziedzin.
Dla $x\in D(f(E))$ ograniczoność $g(E)$ pozwala przejść do granicy
także w $g(E)f_n(E)x$, dając równość podaną w tezie.
W szczególności dla funkcji $f_1,f_2$ i $x$ należącego do obu
dziedzin, rachunek jest addytywny na tym wspólnym zbiorze:
stosujemy wspólne obcięcia $\{|f_1|\le n,\ |f_2|\le n\}$ i zbieżność
zmajoryzowaną przez $2(|f_1|^2+|f_2|^2)$.
Nie oznacza to na ogół równości dziedzin sumy operatorów
i operatora odpowiadającego sumie funkcji.
\end{proof}

\begin{twierdzenie}[Twierdzenie spektralne dla operatora samosprzężonego]
\label{thm:unbounded-spectral33}
Niech $A:D(A)\subseteq H\to H$ będzie gęsto określonym
operatorem samosprzężonym. Istnieje dokładnie jedna miara spektralna
$E:\mathcal B(\R)\to\mathcal L(H)$ taka, że
\[
 D(A)=\left\{x\in H:\int_\R t^2\,\mathrm d\mu_x(t)<\infty\right\},
 \qquad Ax=\int_\R t\,\mathrm dE(t)\,x\quad(x\in D(A)).
\]
Dla każdego $z\in\C\setminus\R$ zachodzi
$(A-zI)^{-1}=\int_\R(t-z)^{-1}\,\mathrm dE(t)$.
Dla dowolnej borelowskiej funkcji $f:\R\to\C$
definiujemy $f(A)$ według
lematu~\ref{lem:unbounded-integral33}; operator ten jest domknięty,
a jego dziedzina jest wyznaczona przez całkę $\int|f|^2\,\mathrm d\mu_x$.
\end{twierdzenie}
\begin{proof}
\textbf{Krok 1: transformacja Cayleya.}
Z twierdzenia~\ref{thm:selfadjoint-range33} operatory
$A\pm iI$ odwzorowują bijektywnie $D(A)$ na $H$.
Definiujemy ograniczony operator
$U=(A-iI)(A+iI)^{-1}=I_H-2i(A+iI)^{-1}$.
Dla $y=(A+iI)x$ symetria daje
$\|Uy\|=\|(A-iI)x\|=\|(A+iI)x\|=\|y\|$.
Surjektywność $A-iI$ pokazuje, że $U$ jest unitarny.
Ponadto $I_H-U=2i(A+iI)^{-1}$ jest różnowartościowy,
a zatem $\ker(I_H-U)=\{0\}$.

\textbf{Krok 2: przeniesienie miary spektralnej.}
Niech $P$ będzie jedyną miarą spektralną operatora
unitarnego $U$ na $\sigma(U)\subseteq\mathbb T$.
Przedłużamy ją na okrąg jednostkowy, kładąc
$P(C):=P(C\cap\sigma(U))$ dla borelowskich $C\subseteq\mathbb T$.
Ponieważ $\ker(I_H-U)=P(\{1\})H=\{0\}$,
mamy $P(\{1\})=0$. Definiujemy rzeczywistą funkcję borelowską
$h(w)=i(1+w)/(1-w)$ dla $w\ne1$ na okręgu jednostkowym.
Jej wartości są rzeczywiste: dla $w=e^{i\theta}\ne1$
liczba $h(w)=-\cot(\theta/2)\in\R$.
Wartość $h(1)$ ustalamy równą $0$, ponieważ $P(\{1\})=0$.
Kładziemy $E(B)=P(h^{-1}(B))$ dla $B\in\mathcal B(\R)$.
Otrzymujemy miarę spektralną na $\R$.

\textbf{Krok 3: identyfikacja operatora.}
Operator $B=h(P)$ z lematu~\ref{lem:unbounded-integral33}
jest samosprzężony. Dla $w\ne1$ zachodzi
$(h(w)+i)^{-1}=(1-w)/(2i)$.
Oznaczmy tę ograniczoną funkcję po prawej przez $g$.
Ponieważ $hg=1-ig$ poza zbiorem $P$-zerowym, iloczyn $hg$
jest ograniczony $P$-prawie wszędzie. Z dowodu lematu
\ref{lem:unbounded-integral33} wynika $g(P)H\subseteq D(B)$
i $(B+iI)g(P)=I_H$, natomiast dla $x\in D(B)$ mamy
$g(P)(B+iI)x=x$. Rachunek spektralny daje więc obustronną odwrotność
$(B+iI)^{-1}=(I_H-U)/(2i)=(A+iI)^{-1}$.
Odwrotności wyznaczają jednoznacznie dziedziny i działania
operatorów, więc $A=B$. Wzór na dziedzinę i całkę
spektralną wynika z lematu i definicji miary obrazowej:
$\mu_x^E(C)=\mu_x^P(h^{-1}(C))$, zatem
$\int_\R t^2\,\mathrm d\mu_x^E=\int_{\mathbb T}|h|^2\,\mathrm d\mu_x^P$.
Dla funkcji prostych całkowanie względem $E$ odpowiada
całkowaniu ich złożeń z $h$ względem $P$; obcięcia i granica
w normie wektora dają identyczność całek nieograniczonych.

\textbf{Krok 4: jednoznaczność.}
Jeżeli $F$ jest inną miarą spektralną reprezentującą $A$,
to ograniczony rachunek borelowski daje
$U=(A-iI)(A+iI)^{-1}=\int_\R (t-i)/(t+i)\,\mathrm dF(t)$.
Jednoznaczność miary spektralnej operatora normalnego
$U$ z \hyperref[sec:spectral-hilbert32]{rozdziału~32}
stosujemy po sprawdzeniu nośnika miary obrazowej $Q$ na $\mathbb T$.
Jeśli $w\notin\sigma(U)$, niech $\delta>0$ będzie takie, że
$\|(U-wI)y\|\ge\delta\|y\|$ dla wszystkich $y\in H$.
Gdyby $Q(\mathbb T\cap B(w,\delta/2))\ne0$, wektor jednostkowy
z obrazu tego rzutu spełniałby $\|(U-wI)y\|\le\delta/2$,
na mocy wzoru całkowego, co jest sprzecznością.
Otwarty zbiór $\mathbb T\setminus\sigma(U)$ jest pokryty
przeliczalnie wieloma takimi otoczeniami (okrąg ma przeliczalną bazę).
Ich suma jest $Q$-zerowa, gdyż wszystkie miary skalarne $\mu_x^Q$
są przeliczalnie addytywne. Zatem $Q(\sigma(U))=I_H$.
Jednoznaczność na $\sigma(U)$ daje $Q=P$ także na $\mathbb T$.
Funkcja ta jest różnowartościowa z $\R$ na
$\mathbb T\setminus\{1\}$, o odwrotności $h$, więc $F=E$.
Dla $z\notin\R$ funkcje $g_z(t)=(t-z)^{-1}$ i $tg_z(t)=1+zg_z(t)$
są ograniczone. Tak jak w kroku 3, $g_z(E)H\subseteq D(A)$,
$(A-zI)g_z(E)=I_H$ i $g_z(E)(A-zI)x=x$ dla $x\in D(A)$.
To dowodzi wzoru na rezolwentę.
\end{proof}

\begin{wniosek}[Widmo i przestrzenie spektralne]
\label{cor:unbounded-spectrum33}
Niech $A$ będzie samosprzężony z miarą spektralną $E$.
Wówczas $\sigma(A)\subseteq\R$ i jest równe nośnikowi $E$,
tj. najmniejszemu domkniętemu zbiorowi $F\subseteq\R$
z $E(F)=I_H$. Dla $\lambda\in\R$ zachodzi
$\ker(A-\lambda I)=E(\{\lambda\})H$.
Jeżeli $\lambda\in\sigma(A)$, to dla każdego $\varepsilon>0$
istnieje $x\in D(A)$ o normie $1$ taki, że
$\|(A-\lambda I)x\|<\varepsilon$.
\end{wniosek}
\begin{proof}
Najpierw istnienie nośnika: suma wszystkich względnie otwartych
zbiorów $E$-zerowych w $\R$ jest sumą przeliczalnie wielu
takich zbiorów, ponieważ $\R$ ma przeliczalną bazę.
Jest więc $E$-zerowa, co wynika z przeliczalnej addytywności
każdej $\mu_x$. Jej dopełnienie jest najmniejszym domkniętym
zbiorem pełnej miary spektralnej. Poza osią rzeczywistą
rezolwenta istnieje z twierdzenia \ref{thm:selfadjoint-range33}.
Dla rzeczywistego $\lambda\notin\operatorname{supp}E$ istnieje $\delta>0$
takie, że $E((\lambda-\delta,\lambda+\delta))=0$.
Kładziemy $g(t)=(t-\lambda)^{-1}$ dla $|t-\lambda|\ge\delta$
oraz $g(t)=0$ w pozostałych punktach. Funkcje $g$ i $tg$
są ograniczone, a $(t-\lambda)g(t)=1$ $E$-prawie wszędzie.
Argument z kroku 3 dowodu twierdzenia \ref{thm:unbounded-spectral33}
daje obustronną odwrotność $g(E)$ operatora $A-\lambda I$. Odwrotnie, jeśli
$\lambda\in\operatorname{supp}E$, to
$E((\lambda-\varepsilon/2,\lambda+\varepsilon/2))\ne0$.
Wybieramy jednostkowy wektor $x$ z obrazu tego rzutu;
jego miara skalarna jest skupiona na tym ograniczonym przedziale,
a więc $x\in D(A)$ i $\|(A-\lambda I)x\|\le\varepsilon/2<\varepsilon$. Operator nie może mieć wtedy
ograniczonej odwrotności. Tożsamość dotycząca jądra wynika
z $\|(A-\lambda I)x\|^2=\int|t-\lambda|^2\,\mathrm d\mu_x(t)$:
zerowa całka oznacza, że $\mu_x$ jest skupiona na $\{\lambda\}$
(rozważamy zbiory $|t-\lambda|\ge1/n$), czyli $E(\{\lambda\})x=x$.
Odwrotnie, każdy wektor z obrazu tego rzutu należy do $D(A)$
i jest zerowany przez $A-\lambda I$.
\end{proof}

\begin{przyklad}[Różne dziedziny tego samego wzoru]
Niech $H=L^2(\R,\lambda_1)$.
Operator $Af=xf$ z dziedziną $D(A)=\{f:xf\in L^2\}$ jest
samosprzężony; jego miara spektralna to
$(E(C)f)(x)=\chi_C(x)f(x)$ (aksjomaty i reprezentacja całkowa
wynikają ze zbieżności zmajoryzowanej i obcięć). Każdy niepusty
otwarty przedział ma dodatnią miarę Lebesgue'a, więc nośnikiem
$E$ jest $\R$ i $\sigma(A)=\R$. Dla operatora
$Bf=-if'$ z dziedziną $H^1(\R)=\{f\in L^2(\R):f'\in L^2(\R)\}$
(pochodna rozumiana w sensie dystrybucyjnym) transformacja Fouriera
z \hyperref[sec:temperowane22]{rozdziału~22} daje w sensie
dystrybucyjnym $\widehat{Bf}(\xi)=\xi\widehat f(\xi)$.
Twierdzenie Plancherela (Twierdzenie~\ref{thm:plancherel22})
i odwrotna transformacja dają równoważność
$f'\in L^2$ oraz $\xi\widehat f\in L^2$; zatem unitarna
transformacja Fouriera przenosi dokładnie $D(B)$ na $D(M_\xi)$.
Wobec tego $B$ jest
samosprzężony i ma widmo $\R$. Tymczasem ograniczenie
operatora $-if'$ do przestrzeni $C_c^\infty(\R)$
jest tylko symetryczne i domykalne, a nie domknięte.
Symetria wynika z całkowania przez części, a gęstość z gęstości
$C_c^\infty(\R)$ w $L^2$. Dla $f(x)=e^{-x^2}$ wybieramy
$\eta\in C_c^\infty(\R)$ równą $1$ w otoczeniu zera i
$f_n(x)=\eta(x/n)f(x)$. Wtedy $f_n\to f$ i
$f_n'=\eta(\cdot/n)f'+n^{-1}\eta'(\cdot/n)f\to f'$ w $L^2$
przez zbieżność zmajoryzowaną oraz oszacowanie drugiego składnika.
Jednak $f$ nie ma zwartego nośnika, więc wykres ograniczenia
nie jest domknięty.
\end{przyklad}

\subsection{Rozszerzenia samosprzężone i indeksy defektu}

\begin{definicja}[Przestrzenie defektu]
\label{def:deficiency33}
Dla gęsto określonego domkniętego operatora symetrycznego
$S$ definiujemy przestrzenie defektu
$N_+=\ker(S^*-iI)$ i $N_-=\ker(S^*+iI)$
oraz indeksy defektu $n_\pm=\dim N_\pm$ rozumiane jako liczby
kardynalne, tj. moce baz ortonormalnych (w szczególności
dopuszczamy wymiary nieprzeliczalne).
\end{definicja}

\begin{twierdzenie}[Rozkład von Neumanna i rozszerzenia]
\label{thm:von-neumann33}
Niech $S$ będzie gęsto określonym domkniętym operatorem
symetrycznym w zespolonej przestrzeni Hilberta $H$.
W normie wykresowej $D(S^*)$ zachodzi ortogonalny rozkład
\[
 D(S^*)=D(S)\oplus N_+\oplus N_-.
\]
Operator $S$ jest samosprzężony wtedy i tylko wtedy,
gdy $N_+=N_-=\{0\}$. Istnieje rozszerzenie samosprzężone
$S\subseteq\widetilde S=\widetilde S^*$ wtedy i tylko wtedy,
gdy $\dim N_+=\dim N_-$. Wszystkie takie rozszerzenia są
parametryzowane operatorami unitarnymi $V:N_+\to N_-$:
\[
 D(S_V)=D(S)\dotplus\{u+Vu:u\in N_+\},\qquad
 S_V(x+u+Vu)=Sx+iu-iVu.
\]
\end{twierdzenie}
\begin{proof}
Operator $S^*$ jest domknięty, a $D(S)$ jest domkniętą
podprzestrzenią przestrzeni Hilberta $D(S^*)$ z normą wykresową.
Dla $x\in D(S)$ i $u\in N_\pm$ otrzymujemy, korzystając
z $S\subseteq S^*$,
$\langle x,u\rangle_{S^*}=\langle x,u\rangle+
\langle Sx,\pm i u\rangle=0$.
Również $N_+\perp N_-$ w iloczynie wykresowym, ponieważ
$\langle iu,-ip\rangle=-\langle u,p\rangle$.
Każda przestrzeń $N_\pm$ jest domknięta w $H$: jeśli
$u_n\in N_\pm$ i $u_n\to u$, to $S^*u_n=\pm iu_n\to\pm iu$,
więc z domkniętości $S^*$ wynika $u\in N_\pm$.
Na $N_\pm$ norma wykresowa jest równa $\sqrt2$ razy normie $H$,
co uzasadnia użycie zupełności tych przestrzeni poniżej.

Niech $z\in D(S^*)$ będzie prostopadły w normie wykresowej
do $D(S)$. To znaczy
$\langle x,z\rangle+\langle Sx,S^*z\rangle=0$
dla każdego $x\in D(S)$. Z definicji operatora sprzężonego
wynika $S^*z\in D(S^*)$ oraz $S^{*2}z=-z$.
Stąd $z_\pm=(z\mp iS^*z)/2$ należą odpowiednio do $N_\pm$
i $z=z_++z_-$. Otrzymujemy cały rozkład.

Dla wektorów $u,v\in N_+$ i $p,q\in N_-$
forma brzegowa
$\mathfrak b(w,z)=\langle S^*w,z\rangle-
\langle w,S^*z\rangle$ spełnia
$\mathfrak b(u+p,v+q)=2i(\langle u,v\rangle-\langle p,q\rangle)$.
Forma $\mathfrak b$ znika, gdy jeden z argumentów należy do
$D(S)$, na mocy definicji sprzężenia. Dla argumentów z $N_+$ i
$N_-$ wyrazy mieszane znoszą się, ponieważ sprzężenie skalaru
$\pm i$ w drugim argumencie zmienia jego znak.
Każde domknięte rozszerzenie $S\subseteq T\subseteq S^*$
odpowiada domkniętej podprzestrzeni
$L\subseteq N_+\oplus N_-$ przez
$D(T)=D(S)\oplus L$. Operator jest symetryczny dokładnie
wtedy, gdy $\mathfrak b$ znika na $L\times L$.
Ponieważ $S\subseteq T$, mamy $T^*\subseteq S^*$.
Dla $z\in D(S^*)$ warunek $z\in D(T^*)$ jest równoważny
$\langle Tw,z\rangle=\langle w,S^*z\rangle$ dla wszystkich
$w\in D(T)$, czyli $\mathfrak b(w,z)=0$ dla wszystkich $w\in L$.
Zatem $D(T^*)=D(S)\oplus L^{\perp_{\mathfrak b}}$, gdzie
$L^{\perp_{\mathfrak b}}=\{z\in N_+\oplus N_-:\mathfrak b(w,z)=0
\ \text{dla wszystkich }w\in L\}$.
Samosprzężoność oznacza więc $L=L^{\perp_{\mathfrak b}}$.

Niech najpierw $L$ odpowiada operatorowi samosprzężonemu.
Jeżeli $u+p\in L$, $u\in N_+$, $p\in N_-$, to
$\mathfrak b(u+p,u+p)=0$ daje $\|u\|=\|p\|$.
Rzut $L$ na $N_+$ jest więc różnowartościowy, a $L$
ma postać wykresu pewnej izometrii
$V:D(V)\subseteq N_+\to N_-$. Podprzestrzeń $D(V)$
jest domknięta: jeśli $u_n\to u$ w $N_+$, to
$Vu_n$ jest ciągiem Cauchy'ego i zbiega do pewnego $p\in N_-$.
Ponieważ $L$ jest domknięta w $N_+\oplus N_-$,
$u+p\in L$, a więc $u\in D(V)$ i $Vu=p$. Jeśli $D(V)\ne N_+$,
istniałby $0\ne w\in D(V)^\perp\subset N_+$.
Wówczas $\mathfrak b(w,l)=0$ dla każdego $l\in L$, ale
$w\notin L$, wbrew $L=L^{\perp_{\mathfrak b}}$.
A więc $D(V)=N_+$. Analogicznie obraz $V$ jest domknięty,
a gdyby nie był całym $N_-$, wektor
$0\ne q\in (\operatorname{Ran}V)^\perp$ należałby do
$L^{\perp_{\mathfrak b}}\setminus L$. Zatem $V$ jest unitarny.
Odwrotnie, dla dowolnego unitarnego $V:N_+\to N_-$
jego wykres $L=\{u+Vu:u\in N_+\}$ jest domknięty
i spełnia $L=L^{\perp_{\mathfrak b}}$: dla $v+q\in N_+\oplus N_-$
warunek $\mathfrak b(u+Vu,v+q)=0$ dla każdego $u\in N_+$
oznacza $\langle u,v-V^*q\rangle=0$, a więc $v=V^*q$,
równoważnie $q=Vv$. Otrzymujemy samosprzężone rozszerzenie
o wskazanej dziedzinie. Operator unitarny między $N_+$ i $N_-$
istnieje dokładnie wtedy, gdy ich bazy ortonormalne mają równe moce:
w jednym kierunku obraz bazy przez unitarny operator jest bazą,
w drugim bijekcja baz wyznacza izometrię na ich liniowych
rozpiętościach, przedłużającą się na całe przestrzenie.
Gdy $N_+=N_- =\{0\}$, rozkład dziedziny daje
$D(S^*)=D(S)$, a więc $S=S^*$.
\end{proof}

\begin{przyklad}[Pochodna na przedziale i warunki brzegowe]
\label{ex:momentum-interval33}
Na $H=L^2(0,1)$ rozważmy
$Sf=-if'$ z dziedziną
$D(S)=\{f\in H^1(0,1):f(0)=f(1)=0\}$.
Przestrzeń $H^1(0,1)$ oznacza funkcje z $L^2(0,1)$
o słabej pochodnej w $L^2$, identyfikowane z ich
absolutnie ciągłymi reprezentantami.
Gęstość dziedziny wynika z zawierania $C_c^\infty(0,1)$.
Całkowanie przez części daje symetrię. Norma wykresowa jest
normą $H^1$, która jest zupełna: jeżeli $f_n\to f$ i $f_n'\to g$
w $L^2$, testowanie na funkcjach testowych daje $f'=g$ w sensie
dystrybucyjnym. Wartości brzegowe są ciągłe w tej normie, gdyż
dla $t\in[0,1]$ mamy
$|f(0)|\le |f(t)|+\int_0^1|f'(s)|\,\mathrm d\lambda_1(s)$
i analogicznie dla $f(1)$. Całkując względem $t$ oraz stosując
nierówność Cauchy'ego--Schwarza, dostajemy
$|f(0)|,|f(1)|\le\|f\|_{L^2}+\|f'\|_{L^2}$.
Zatem $D(S)$ jest domkniętą podprzestrzenią $H^1$, a operator
$S$ jest domknięty na mocy kryterium normy wykresowej.
Operator sprzężony ma ten sam wzór $S^*f=-if'$, lecz
$D(S^*)=H^1(0,1)$. Istotnie, testowanie warunku
sprzężenia na funkcjach z $C_c^\infty(0,1)$ wymusza istnienie
słabej pochodnej w $L^2$, a zatem przynależność do $H^1$.
Odwrotnie, całkowanie przez części dla dwóch funkcji z $H^1$
pokazuje, że wyrazy brzegowe znikają, ponieważ
wektor z $D(S)$ ma zerowe wartości na obu końcach.
Równania $S^*f=if$ i $S^*f=-if$ mają odpowiednio
rozwiązania $f(x)=ce^{-x}$ i $f(x)=ce^x$.
Zatem $n_+=n_-=1$. Wzór na formę brzegową brzmi
$\mathfrak b(f,g)=-i(f(1)\overline{g(1)}-f(0)\overline{g(0)})$.
Rozszerzenia samosprzężone są dokładnie operatorami
$-if'$ o dziedzinach
$D(S_\theta)=\{f\in H^1(0,1):f(1)=e^{i\theta}f(0)\}$,
$\theta\in[0,2\pi)$. Aby to uzasadnić, odwzorowanie śladu
$\gamma:H^1(0,1)\to\C^2$, $\gamma(f)=(f(0),f(1))$,
jest surjektywne (wystarczą funkcje afiniczne) i ma jądro $D(S)$.
Dziedziny rozszerzeń samosprzężonych odpowiadają więc podprzestrzeniom
$L\subseteq\C^2$ równym własnemu anihilatorowi względem formy
$-i(b\overline d-a\overline c)$ na parach $(a,b),(c,d)$.
Taka podprzestrzeń ma wymiar $1$, gdyż forma jest niezdegenerowana.
Dla jej niezerowego generatora $(a,b)$ izotropowość daje $|a|=|b|$,
więc $a\ne0$ i $b=e^{i\theta}a$. Odwrotnie każda taka prosta
jest izotropowa i ma wymiar $1$, stąd jest równa własnemu
anihilatorowi. Są to różne realizacje \emph{tego samego wzoru}.
\end{przyklad}

\begin{przyklad}[Brak rozszerzenia samosprzężonego]
Na $L^2(0,\infty)$ operator $-if'$ z dziedziną
$H_0^1(0,\infty)=\{f\in H^1(0,\infty):f(0)=0\}$
jest domknięty i symetryczny. Przyjmujemy, że $H^1(0,\infty)$
składa się z funkcji $L^2$ o słabej pochodnej w $L^2$;
każda z nich ma lokalnie absolutnie ciągłego reprezentanta.
Dziedzina jest gęsta, gdyż zawiera $C_c^\infty(0,\infty)$.
Zupełność $H^1$ wynika z tego samego argumentu dystrybucyjnego
co na odcinku, a ślad w zerze jest ciągły dzięki oszacowaniu
na $[0,1]$. To daje domkniętość. Dla $f\in H^1(0,\infty)$
pochodna $|f|^2$ jest równa $2\operatorname{Re}(f'\overline f)$
i należy do $L^1$ przez nierówność Cauchy'ego--Schwarza.
Zatem $|f(t)|^2$ ma skończoną granicę przy $t\to\infty$,
a z całkowalności $|f|^2$ wynika, że granica ta jest zerowa.
Całkowanie przez części na $[0,R]$ i przejście z $R$ do
nieskończoności usuwa więc wyraz brzegowy w nieskończoności.
Dla dwóch funkcji z dziedziny znika również wyraz w zerze,
co daje symetrię. Testowanie sprzężenia na funkcjach testowych
wymusza słabą pochodną w $L^2$, a całkowanie przez części
z dowolną funkcją z $H^1$ daje inkluzję odwrotną.
Jego operator sprzężony ma zatem dziedzinę $H^1(0,\infty)$. Z równań defektu otrzymujemy
$N_+=\lin\{e^{-x}\}$ i $N_- =\{0\}$, ponieważ $e^x\notin L^2$.
Indeksy defektu różnią się, więc operator nie ma rozszerzenia
samosprzężonego w tej samej przestrzeni Hilberta.
\end{przyklad}

\subsection{Jednoparametrowe grupy unitarne i twierdzenie Stone'a}

\begin{definicja}[Silnie ciągła grupa unitarna]
Niech $(U_t)_{t\in\R}$ będzie rodziną operatorów unitarnych
na $H$ z $U_0=I_H$ i $U_{t+s}=U_tU_s$ dla $s,t\in\R$.
Jest to \emph{silnie ciągła jednoparametrowa grupa unitarna}, gdy
$\lim_{t\to0}\|U_tx-x\|=0$ dla każdego $x\in H$.
Jej generator definiujemy na dziedzinie
$D(G)=\{x:\lim_{t\to0}(U_tx-x)/t\text{ istnieje w }H\}$
przez $Gx=\lim_{t\to0}(U_tx-x)/t$.
\end{definicja}

\begin{twierdzenie}[Twierdzenie Stone'a]
\label{thm:stone33}
Dla każdej silnie ciągłej jednoparametrowej grupy unitarnej
$(U_t)_{t\in\R}$ na zespolonej przestrzeni Hilberta $H$
istnieje dokładnie jeden operator samosprzężony $A$ taki, że
$U_t=e^{itA}$ dla wszystkich $t\in\R$, gdzie
$e^{itA}=\int_\R e^{it\lambda}\,\mathrm dE_A(\lambda)$.
Ponadto $D(A)=D(G)$ i $G=iA$.
Odwrotnie, dla każdego samosprzężonego $A$ rachunek spektralny
wyznacza silnie ciągłą grupę unitarną $(e^{itA})_{t\in\R}$,
a jej generator jest równy $iA$.
\end{twierdzenie}
\begin{proof}
\textbf{Krok 1: konstrukcja grupy z operatora.}
Dla samosprzężonego $A$ funkcja $\lambda\mapsto e^{it\lambda}$
jest ograniczona i ma moduł $1$, zatem rachunek borelowski
definiuje operator $U_t=e^{itA}$ spełniający
$U_t^*U_t=U_tU_t^*=1(E_A)=I_H$. Operator jest więc unitarny.
Równości $U_{t+s}=U_tU_s$ i $U_0=I_H$ są punktowymi
tożsamościami funkcji w rachunku spektralnym.
Dla każdego $x\in H$ mamy
$\|(U_t-I_H)x\|^2=\int_\R|e^{it\lambda}-1|^2\,\mathrm d\mu_x(\lambda)
\to0$ na mocy zbieżności zmajoryzowanej.
Dla $x\in D(A)$ nierówność
$|(e^{it\lambda}-1)/t|\le|\lambda|$ (z całkowania po odcinku
funkcji $s\mapsto i\lambda e^{is\lambda}$) i zbieżność punktowa
ilorazu różnicowego do $i\lambda$ dają
\[
 \left\|\frac{U_tx-x}{t}-iAx\right\|^2
 =\int_\R\left|\frac{e^{it\lambda}-1}{t}-i\lambda\right|^2
 \,\mathrm d\mu_x(\lambda)\longrightarrow0.
\]
Równość wynika z addytywności na wspólnej dziedzinie
w dowodzie Lematu~\ref{lem:unbounded-integral33}, a funkcja pod
całką jest ograniczona przez całkowalną funkcję $4\lambda^2$.
Odwrotnie, gdy dla $x$ iloraz ma granicę, lemat Fatou
dla miary $\mu_x$ i
$\lim_{t\to0}|(e^{it\lambda}-1)/t|^2=\lambda^2$
dają, biorąc np. ciąg $t=1/n$,
$\int_\R\lambda^2\,\mathrm d\mu_x\le
\liminf_n\|n(U_{1/n}x-x)\|^2<\infty$.
Zatem $D(G)=D(A)$.

\textbf{Krok 2: gęstość i skośna symetria generatora.}
Niech teraz $(U_t)$ będzie dowolną silnie ciągłą grupą unitarną.
Dla $\phi\in C_c^\infty(\R)$ i $x\in H$ definiujemy
$x_\phi=\int_\R\phi(t)U_tx\,\mathrm d\lambda_1(t)$.
Jest to całka w normie przestrzeni Hilberta, zdefiniowana
przez granice sum Riemanna na zwartym nośniku $\phi$;
nie potrzeba tu ogólnej teorii całki Pettisa. Ciągłość
$t\mapsto U_tx$ w każdym punkcie wynika z
$\|U_{t+s}x-U_tx\|=\|U_sx-x\|$.
Przesunięcie zmiennej pod całką daje dla $s\ne0$
\[
 \frac{U_sx_\phi-x_\phi}{s}
 =\int_\R\frac{\phi(t-s)-\phi(t)}s U_tx\,\mathrm d\lambda_1(t).
\]
Dla $|s|\le1$ ilorazy mają wspólny zwarty nośnik i zbiegają
jednostajnie do $-\phi'$ (z jednostajnej ciągłości $\phi'$).
Oszacowanie normy całki przez całkę norm dowodzi zbieżności
do $-\int_\R\phi'(t)U_tx\,\mathrm d\lambda_1(t)$ w $H$.
W szczególności $x_\phi\in D(G)$. Wybierając
nieujemną aproksymację jedności $\phi_n$ wokół zera,
otrzymujemy $x_{\phi_n}\to x$ z silnej ciągłości grupy.
Dlatego $D(G)$ jest gęsta.
Dla $x,y\in D(G)$ różniczkowanie
$\langle U_tx,U_ty\rangle=\langle x,y\rangle$ w zerze daje
$\langle Gx,y\rangle+\langle x,Gy\rangle=0$.
Stąd operator $A=-iG$ jest symetryczny.

\textbf{Krok 3: surjektywność $A\pm iI$.}
Dla $x\in H$ kładziemy
$R_+x=\int_0^\infty e^{-t}U_tx\,\mathrm d\lambda_1(t)$,
$R_-x=\int_0^\infty e^{-t}U_{-t}x\,\mathrm d\lambda_1(t)$.
Są to granice w normie całek na odcinkach $[0,N]$;
$\|R_\pm\|\le1$, ponieważ $\|U_tx\|=\|x\|$.
Dla $s\in\R$ zmiana zmiennej daje
$U_sR_+x=e^s\int_s^\infty e^{-r}U_rx\,\mathrm d\lambda_1(r)$.
Iloraz różnicowy w zerze zbiega więc w normie do
$R_+x-x$ (z silnej ciągłości $U_rx$ w $r=0$).
Stąd $R_+x\in D(G)$ i $(I_H-G)R_+x=x$.
Podobnie
$U_sR_-x=e^{-s}\int_{-s}^\infty e^{-r}U_{-r}x\,\mathrm d\lambda_1(r)$,
zatem $GR_-x=x-R_-x$ i $(I_H+G)R_-x=x$.
W konsekwencji
$\operatorname{Ran}(I_H\pm G)=H$. Ponieważ $A=-iG$,
mamy $A+iI=i(I_H-G)$ i $A-iI=-i(I_H+G)$ na $D(G)$,
więc $\operatorname{Ran}(A\pm iI)=H$.
Twierdzenie~\ref{thm:selfadjoint-range33} daje samosprzężoność $A$.

\textbf{Krok 4: identyfikacja i jednoznaczność grupy.}
Z własności grupowej wynika
$U_tD(G)\subseteq D(G)$ i $GU_tx=U_tGx$:
istotnie, iloraz różnicowy w punkcie $t$ jest obrazem
ilorazu różnicowego w zerze przez ograniczony operator $U_t$.
Analogiczna własność zachodzi dla $e^{itA}$
na mocy rachunku spektralnego.
Dla $x\in D(A)$ zarówno $t\mapsto U_tx$, jak i
$t\mapsto e^{itA}x$ są więc różniczkowalne i spełniają
równanie $v'(t)=iAv(t)$ oraz warunek $v(0)=x$.
Ich różnica $w$ ma stałą normę, ponieważ
$\frac{\mathrm d}{\mathrm dt}\|w(t)\|^2
=2\operatorname{Re}\langle iAw(t),w(t)\rangle=0$.
Ponieważ $w(0)=0$, mamy $w=0$.
Gęstość $D(A)$ i unitarność obu grup rozszerzają równość
na całe $H$. Jednoznaczność $A$ wynika z definicji generatora.
\end{proof}
\endgroup
